% Lösungen Einheit 1 von 2 – Jahreszins, Monats- und Tageszins (zusammenbau v0.1)
\einheitenkopf{Einheit 1 von 2 · Jahreszins, Monats- und Tageszins}

%% TODO Lösung der Erklärzeile 11g) fehlt in der Bank
\erg{11}{a) die Zinsen (Prozentwert) \quad b) 6\,€ ($300 \cdot 0{,}02$) \quad c) 21\,€ ($700 \cdot 0{,}03$) \quad d) 80\,€ ($2\,000 \cdot 0{,}04$) \quad e) 18\,€ ($900 \cdot 0{,}02$) \quad f) 150\,€ ($5\,000 \cdot 0{,}03$) \quad g) \ldots \quad h) 24\,€ (1\,\% = 6\,€, 4\,\% = 24\,€) \quad i) 714\,€ (Zinsen: $700 \cdot 0{,}02 = 14$\,€) \quad j) 3\,\% ($27 : 900 = 0{,}03$) \quad k) 1,5\,\% ($60 : 4\,000 = 0{,}015$) \quad l) 700\,€ (1\,\% = 7\,€) \quad m) $300 \cdot 0{,}03 = 9{,}00$\,€ – stimmt (oder $9 : 300 = 0{,}03 = 3\,\%$). \quad n) 36\,€ (Jahreszinsen 12\,€, $12 \cdot 3 = 36$) \quad o) 1\,908\,€ (Zinsen: $1\,800 \cdot 0{,}06 = 108$\,€) \quad p) Bank B: 65\,€, Bank A: 60\,€. \quad q) 52\,€ ($2\,600 \cdot 0{,}02$)}
%% TODO Lösung der Erklärzeile 12g) fehlt in der Bank
\erg{12}{a) Monate – durch zwölf teilen \quad b) 20\,€ (Jahreszinsen 48\,€, $48 : 12 \cdot 5 = 20$) \quad c) 12\,€ (Jahreszinsen 36\,€, $36 : 12 \cdot 4 = 12$) \quad d) 42\,€ (Jahreszinsen 72\,€, $72 : 12 \cdot 7 = 42$) \quad e) 60\,€ (Jahreszinsen 240\,€, $240 : 12 \cdot 3 = 60$) \quad f) 108\,€ (Jahreszinsen 144\,€, $144 : 12 \cdot 9 = 108$) \quad g) \ldots \quad h) 16\,€ (Jahreszinsen 144\,€, $144 : 360 \cdot 40 = 16$) \quad i) 45\,€ (75 Tage; Jahreszinsen 216\,€, $216 : 360 \cdot 75 = 45$) \quad j) 2\,\% (Jahreszinsen $20 : 4 \cdot 12 = 60$\,€, $60 : 3\,000 = 0{,}02$) \quad k) 52,50\,€ (Jahreszinsen 105\,€, $105 : 12 \cdot 6 = 52{,}50$). Nora bekommt 52,50\,€ Zinsen.}
\erg{13}{a) ≈ 80\,€ (1\,\% ≈ 40\,€, mal 2)}
\erg{14}{a) Gefragt sind die Zinsen, nicht das Guthaben: richtig 38\,€. \quad b) Schon 1\,\% von 20\,000\,€ sind 200\,€, 4\,\% von 600\,€ nur 24\,€ – der kleine Anteil vom viel größeren Kapital ist mehr. \quad c) Nein: Rückzahlung 4\,240\,€, gespart 4\,200\,€, es fehlen 40\,€.}
